现在陈述一个关于算子的伴随与其闭包之间关系的定理.

\begin{Theorem}
\label{thm:adjoint-closedness-and-double-adjoint-closure}
设 $(\mathcal{D}(A),A)$ 为从 $X$ 到 $Y$ 的稠密定义域算子. 那么伴随算子 $A^\times$ 是闭的. 此外, 若 $X$ 和 $Y$ 是自反 Banach 空间, 则 $(\mathcal{D}(A),A)$ 可闭, 当且仅当 $\mathcal{D}(A^\times)$ 稠密, 此时 $\overline A=A^{\times\times}$.
\end{Theorem}

\begingroup
\postdisplaypenalty=10000
\renewcommand{\endProof}{\par\nobreak\hfill$\blacksquare$\par\vspace{0.1cm}}
\begin{Proof}
为证明 $A^\times$ 是闭的, 使用 Lemma~\ref{lem:closed-operator-sequential-criterion}. 即假设 $\mathcal{D}(A^\times)\subset Y^*$ 中存在序列 $v_n$, 使得对某个 $v\in Y^*$ 有 $v_n\to v$, 且对某个 $f\in X^*$ 有 $A^\times v_n\to f$. 需要证明, 对每个 $u\in\mathcal{D}(A)$ 均有 $v(Au)=f(u)$. 但由于 $Y^*$ $(X^*)$ 中的收敛蕴含 $Y^*$ $(X^*)$ 中的弱-* 收敛, 对任意 $u\in\mathcal{D}(A)$ 均有\[v(Au)=\lim_{n\to\infty}v_n(Au)=\lim_{n\to\infty}A^\times v_n(u)=f(u)\]因而 $A^\times$ 是闭的.

证明的其余部分留给读者 (Problem~\ref{ex:adjoint-closure-theorem-proof-completion}).
\end{Proof}
\endgroup

\subsubsection{Hilbert 伴随算子}\label{subsec:hilbert-adjoint-operator}

本节只考察一个 Hilbert 空间到另一个 Hilbert 空间的算子. 在这种情形下, 定义以下与伴随密切相关的算子.

\begin{Definition}
\label{def:hilbert-adjoint-operator}
设 $(\mathcal{D}(A),A)$ 为从 Hilbert 空间 $H_1$ 到 Hilbert 空间 $H_2$ 的稠密定义算子. 将 $\mathcal{D}(A^*)$ 定义为满足以下条件的所有 $v\in H_2$ 所组成的集合: 存在 $w\in H_1$, 使得对所有 $u\in\mathcal{D}(A)$ 均有
\begin{equation}\label{eq:hilbert-adjoint-pairing-definition}
(v,Au)_{H_2}=(w,u)_{H_1}
\end{equation}
注意, 由于 $\mathcal{D}(A)$ 稠密, $w$ 由 $v\in\mathcal{D}(A^*)$ 和 \eqref{eq:hilbert-adjoint-pairing-definition} 唯一确定. 因而可用下式定义从 $H_2$ 到 $H_1$ 的算子 $(\mathcal{D}(A^*),A^*)$:
\begin{equation}\label{eq:hilbert-adjoint-action-definition}
A^*v:=w
\end{equation}
其中 $v\in\mathcal{D}(A^*)$. 称 $(\mathcal{D}(A^*),A^*)$ 为 $(\mathcal{D}(A),A)$ 的 \textbf{Hilbert 伴随}.
\end{Definition}

利用 Example~\ref{ex:riesz-map-conjugate-linear-operator} 中定义的 Riesz 映射 $A_{H_1}$ 和 $A_{H_2}$, 容易得到伴随与 Hilbert 伴随之间的关系:
\begin{equation}\label{eq:hilbert-and-banach-adjoint-riesz-relation}
A^*=A_{H_1}\circ A^\times\circ A_{H_2}^{-1}.
\end{equation}
(这里 $\circ$ 表示算子的复合.) 由于 Hilbert 空间与其对偶空间等距, 很少直接使用对偶空间; 在研究 Hilbert 空间上的算子时, 几乎总是使用 Hilbert 伴随而非伴随. 事实上, 通常将 Hilbert 伴随简称为伴随, 完全忽略二者的区别. 这里也采用这一约定, 不过对这两类算子使用不同的记号.

\begin{Theorem}
\label{thm:range-complement-adjoint-nullspace}
对于从 $H$ 到 $H$ 的任意稠密定义算子 $(\mathcal{D}(A),A)$, 值域的正交补为伴随的零空间:
\begin{equation}\label{eq:range-orthogonal-complement-adjoint-nullspace}
\mathcal{R}(A)^\perp=\mathcal{N}(A^*).
\end{equation}
此外, 若 $\mathcal{R}(A)$ 是闭的, 则
\begin{equation}\label{eq:closed-range-adjoint-nullspace-complement}
\mathcal{R}(A)=\mathcal{N}(A^*)^\perp,
\end{equation}
即方程 $Ax=f$ 有解 $x$, 当且仅当 $f\in\mathcal{N}(A^*)^\perp$.
\end{Theorem}

\begingroup
\postdisplaypenalty=10000
\renewcommand{\endProof}{\par\nobreak\hfill$\blacksquare$\par\vspace{0.1cm}}
\begin{Proof}
先证明 $\mathcal{R}(A)^\perp\subset\mathcal{N}(A^*)$. 设 $z\in\mathcal{R}(A)^\perp$. 那么对所有 $u\in\mathcal{D}(A)$, 均有
\begin{equation}\label{eq:range-complement-in-adjoint-domain}
(z,Au)=0=(0,u)
\end{equation}
因而 $z\in\mathcal{D}(A^*)$ 且 $A^*z=0$; 即 $z\in\mathcal{N}(A^*)$. 类似地, 若 $z\in\mathcal{N}(A^*)$, 则对每个 $u\in\mathcal{D}(A)$ 均有
\begin{equation}\label{eq:adjoint-nullspace-in-range-complement}
(z,Au)=(A^*z,u)=(0,u)=0.
\end{equation}
因而 $z\in\mathcal{R}(A)^\perp$, 且 $\mathcal{N}(A^*)\subset\mathcal{R}(A)^\perp$.

最后, 若 $\mathcal{R}(A)$ 是闭的, 则有
\begin{equation}\label{eq:closed-range-double-orthogonal-complement}
\mathcal{R}(A)=\mathcal{R}(A)^{\perp\perp}=\mathcal{N}(A^*)^\perp.
\end{equation}
证明完毕.
\end{Proof}
\endgroup
